\section{On almost disjointness and Urysohn $Q$-spaces}\label{sec:urysohn-q-spaces}
As discussed in the introduction, the cardinal $\mathfrak q$ is the least cardinality of a non-$Q$-set of reals. For $f,g\in\omega^\omega$, write $f\leq^* g$ if $f(n)\leq g(n)$ for all but finitely many $n\in\omega$. It is known that $\mathfrak p\leq \mathfrak q\leq \mathfrak b$, where $\mathfrak b$ is the bounding number, namely the least cardinality of a family in $\omega^\omega$ that is not bounded by any single function with respect to $\leq^*$ \cite{miller1984special}.

In \cite{banakh2023q}, Banakh and Bazylevych introduced the following separation-axiom variants of $\mathfrak q$ (with the classical invariant denoted by $\mathfrak q_0$ in their notation).

\begin{definition}\label{def:q-variants}
    The cardinal $\mathfrak q_1$ is the least cardinal $\kappa$ for which there exists a second-countable, $T_1$, non-$Q$-space of cardinality $\kappa$.

    The cardinal $\mathfrak q_2$ is the least cardinal $\kappa$ for which there exists a second-countable, Hausdorff, non-$Q$-space of cardinality $\kappa$.

    The cardinal $\mathfrak q_{2\frac{1}{2}}$ is the least cardinal $\kappa$ for which there exists a second-countable, Urysohn, non-$Q$-space of cardinality $\kappa$.
\end{definition}

They proved that $\mathfrak{adp}\leq \mathfrak q_1$. It is also immediate from the separation axioms that $\mathfrak q_1\leq \mathfrak q_2\leq \mathfrak q_{2\frac{1}{2}}\leq \mathfrak q$; together with $\mathfrak q\leq\mathfrak b$, this gives, by Theorem~\ref{thm:adp-equals-dp},
\[
\mathfrak p\leq \mathfrak{dp}=\mathfrak{adp}\leq \mathfrak q_1\leq \mathfrak q_2\leq \mathfrak q_{2\frac{1}{2}}\leq \mathfrak q\leq \mathfrak b.
\]
Banakh and Bazylevych asked whether any of these inequalities between different versions of $\mathfrak q$ can be strict.
They also proved that every submetrizable space, that is, every space admitting a coarser metrizable topology, of cardinality less than $\mathfrak q$ is a $Q$-space.
Thus, any separation of these cardinals from $\mathfrak q$ must use non-submetrizable examples \cite{banakh2023q}.

More specifically, they asked whether $\mathfrak{ap}\leq \mathfrak q_2$.
We show in this section that the answer is positive.

The first step is to extract, from the Urysohn separation axiom, a countable base arranged in finite levels.
The levels still refine neighborhoods of each point, but distinct points can occur together in closures only finitely often.

\begin{lemma}\label{lem:urysohn-ad-base}
Let $X$ be a second-countable Urysohn space.
Then $X$ has a countable base $\mathcal B$ and a sequence $(\mathcal B_n)_{n<\omega}$ such that:
\begin{enumerate}[label=(\roman*)]
    \item $\mathcal B=\bigcup_{n<\omega}\mathcal B_n$;
    \item for every $n<\omega$, $\mathcal B_n$ is finite;
    \item for every two distinct points $x,y\in X$, the set $\{n \in \omega: \exists B\in\mathcal B_n\,(x\in\overline B\text{ and }y\in\overline B)\}$ is finite;
    \item for every $x\in X$ and every open neighborhood $O$ of $x$, there is some $m<\omega$ such that for every $n\geq m$ there is some $B\in\mathcal B_n$ such that $x\in B\subseteq O$ (so, in particular, $\mathcal B$ is a base for $X$);
    \item for every $n<\omega$ and $x\in X$, there is some $B\in\mathcal B_n$ such that $x\in B$ (that is, $\mathcal B_n$ covers $X$).
\end{enumerate}
\end{lemma}

\begin{proof}
Since $X$ is second-countable, $X^2\setminus\Delta$ is Lindelöf.
Since $X$ is Urysohn, there exists a sequence of pairs of open subsets of $X$, $(U_k,V_k)_{k<\omega}$,
such that
\[
    X^2\setminus\Delta
    \subseteq \bigcup_{k<\omega} U_k\times V_k \qquad\text{and}\qquad \forall k<\omega\,\overline{U_k}\cap\overline{V_k}=\emptyset.
\]

For each $k<\omega$, let $C_k^0=X\setminus\overline{U_k}$ and $C_k^1=X\setminus\overline{V_k}$.
Since $\overline{U_k}\cap\overline{V_k}=\emptyset$, the pair $\{C_k^0,C_k^1\}$ is an open cover of $X$.

Fix a countable base $(G_m)_{m<\omega}$ for $X$ with $G_0=X$.
For each $n<\omega$, define
\[
    \mathcal B_n=
    \left\{
        G_m\cap \bigcap_{k\leq n} C_k^{\varepsilon(k)}
        :
        m\leq n,\ \varepsilon\in 2^{n+1}
    \right\}.
\]

Each $\mathcal B_n$ is a finite collection of open sets.
Let $\mathcal B=\bigcup_{n<\omega}\mathcal B_n$.
We first verify (iv).
Let $x\in X$, and let
$O\subseteq X$ be open with $x\in O$.
Choose $m<\omega$ such that
$x\in G_m\subseteq O$.
Fix $n\geq m$.
For every $k\leq n$, choose $\varepsilon(k)\in\{0,1\}$ such that
$x\in C_k^{\varepsilon(k)}$. Then
\[
    x\in G_m\cap\bigcap_{k\leq n}C_k^{\varepsilon(k)}
    \subseteq G_m\subseteq O,
\]
and this set belongs to $\mathcal B_n$.
Thus, (iv) holds.
The same argument with $G_0=X$ shows that (v) holds.

It remains to verify (iii).
Let $x\neq y$. Choose
$k<\omega$ such that $(x,y)\in U_k\times V_k$.
Then
\[
    x\notin\overline{C_k^0}
    \quad\text{and}\quad
    y\notin\overline{C_k^1}.
\]
Now let $n>k$ and $B\in\mathcal B_n$. By construction, either
$B\subseteq C_k^0$ or $B\subseteq C_k^1$. In the first case,
$x\notin\overline B$; in the second case, $y\notin\overline B$. Hence no
member of $\mathcal B_n$, for $n>k$, has both $x$ and $y$ in its closure.

Therefore,
\[
    \{n\in\omega:\exists B\in\mathcal B_n\,
    (x\in\overline B\text{ and }y\in\overline B)\}
    \subseteq \{0,\ldots,k\},
\]
which is finite.
\end{proof}

We now define a cardinal $\mathfrak{at}$ that is a tree analogue of $\mathfrak{ap}$.
We will show that $\mathfrak{at}=\mathfrak q_{2}$.

The preceding lemma turns Urysohn separation into a countable combinatorial object.
This motivates the following tree variant of $\mathfrak{ap}$, which keeps track of weak separation for almost disjoint subtrees rather than for almost disjoint subsets of $\mathbb N$.

In this paper, an $\omega$-tree is a rooted tree of height $\omega$ with finite levels, and subtrees are assumed to be nonempty and downward closed.  For $m<\omega$, write
\[
  T_m=\{t\in T:\htt(t)=m\}
  \quad\text{and}\quad
  T_{<m}=\{t\in T:\htt(t)<m\}.
\]
Two subtrees $S,R$ of $T$ are almost disjoint if $|S\cap R|<\omega$.
A family of subtrees of $T$ is almost disjoint if any two distinct members are almost disjoint.
An infinite subtree means a subtree with infinitely many nodes.

Recall that a branch of a tree $T$ is a maximal linearly ordered subset of $T$.

\begin{definition}\label{def:at}
Let $\mathfrak{at}$ be the least cardinal $\kappa$ for which there is an
$\omega$-tree $T$, an almost disjoint family $\mathcal X$ of infinite subtrees of $T$ with $|\mathcal X|\leq \kappa$, and a subcollection
$\mathcal Y\subseteq\mathcal X$ such that $\mathcal Y$ cannot be weakly
separated from $\mathcal X\setminus\mathcal Y$.
Here a set $D\subseteq T$
weakly separates $(\mathcal Y,\mathcal X\setminus\mathcal Y)$ if
\[
  \forall Y\in\mathcal Y\ |D\cap Y|=\omega,
  \qquad \text{and}\qquad
  \forall Z\in\mathcal X\setminus\mathcal Y\ |D\cap Z|<\omega.
\]
\end{definition}

It follows immediately that $\mathfrak{ap}\leq \mathfrak{at}$, since every $\omega$-tree is countably infinite and, in the definition of $\mathfrak{ap}$, $\omega$ can be replaced by any infinite countable set.

\begin{proposition}\label{prop:q2-below-at}
$\mathfrak q_2\leq \mathfrak{at}$.
\end{proposition}
\begin{proof}
Let $\kappa<\mathfrak q_2$.
We show that $\kappa<\mathfrak{at}$.
Notice that $\kappa<\mathfrak q_2\leq\mathfrak q\leq\mathfrak b$.

Let $\mathcal X$ be an almost disjoint family
$\mathcal X$ of infinite subtrees of an $\omega$-tree $T$.
Let $\mathcal Y\subseteq\mathcal X$ and put
$\mathcal Z=\mathcal X\setminus\mathcal Y$.
We will construct a set $D\subseteq T$ that weakly separates $\mathcal Y$ from $\mathcal Z$.

If $\mathcal Y=\emptyset$, take $D=\emptyset$; if
$\mathcal Z=\emptyset$, take $D=T$. Hence we may assume that both
subcollections are nonempty.

For each $Y\in\mathcal Y$, choose an infinite branch
$b_Y\subseteq Y$, using König's lemma, and put
\[
  H=\{b_Y:Y\in\mathcal Y\},\qquad X=H\cup\mathcal Z.
\]
The almost disjointness of $\mathcal X$ implies $b_Y\neq b_{Y'}$ whenever $Y, Y'$ are distinct members of $\mathcal Y$, and that $H\cap \mathcal Z=\emptyset$.
Thus, $X$ is an almost disjoint
family of infinite subtrees of $T$ and $|X|\leq\kappa$.

For each finite $F\subseteq T$, define
\[
  N_F=\{A\in X:A\cap F=\emptyset\}.
\]
Notice that 
$N_\emptyset=X$ and $N_F\cap N_G=N_{F\cup G}$.
Thus, these sets form a countable base for a topology on $X$.

We claim that this topology is Hausdorff.
For each $m \in \omega$, let $T_m$ be the $m$-th level of $T$.

Let $A,B\in X$ be distinct, and choose $m<\omega$ such that
$A\cap B\cap T_m=\emptyset$.
Then $N_{T_m\setminus A}$ and $N_{T_m\setminus B}$ are open
neighborhoods of $A$ and $B$, respectively.
If $C$ belonged to both neighborhoods, we would have
\[
  C\cap T_m\subseteq A\cap B\cap T_m=\emptyset.
\]
This is impossible by König's lemma.

Since $|X|\leq\kappa<\mathfrak q_2$, the space $X$ is a $Q$-space.
In particular, write
\[
  H=\bigcup_{n<\omega}H_n,
\]
where each $H_n$ is closed in $X$, and define
\[
  S_n=\bigcup_{b\in H_n}b.
\]
\begin{claim}
  For every $n<\omega$ and every $Z\in\mathcal Z$, the set $S_n\cap Z$ is finite.
\end{claim}
\begin{proof}[Proof of the claim]
Since $Z\notin H_n$ and $H_n$ is closed, there exists a finite $F\subseteq T$ such that
\[
  Z\in N_F\subseteq X\setminus H_n.
\]

Let $M=\max\{\htt(t):t\in F\}+1$.

Given $t\in S_n\cap Z$, $\htt(t)<M$, for if $\htt(t)\geq M$, choose, as $t\in S_n$, an element $b\in H_n$
with $t\in b$.
As $b \in H_n$, it follows that $b\notin N_F$, so there exists $u\in b\cap F$.

Since $b$ is a branch, since $u, t\in b$, and since $\htt(t)\geq M>\htt(u)$, we have $u<t$.
As $t\in Z$ and $Z$ is a subtree, we have $u\in Z$, contradicting $Z\in N_F$.
Consequently,
\[
  S_n\cap Z\subseteq T_{<M},
\]
which is finite.
\end{proof}

For each $Z\in\mathcal Z$, choose $f_Z\in\omega^\omega$ such that
\[
  S_n\cap Z\subseteq T_{<f_Z(n)}
  \qquad(n<\omega).
\]
As $|\mathcal Z|\leq\kappa<\mathfrak b$, there exists
$g\in\omega^\omega$ such that $f_Z\leq^*g$ for every
$Z\in\mathcal Z$. Define
\[
  D=\bigcup_{n<\omega}\bigl(S_n\setminus T_{<g(n)}\bigr).
\]

Given $Y\in\mathcal Y$, there exists $n<\omega$ such that $b_Y\in H_n$.
Hence $b_Y\setminus T_{<g(n)}\subseteq S_n\setminus T_{<g(n)}$, so $|D\cap Y|=\omega$.

Given $Z\in\mathcal Z$, choose $N<\omega$ such that
$f_Z(n)\leq g(n)$ for all $n\geq N$.
For such numbers $n$, the set $S_n\setminus T_{<g(n)}$ is disjoint from $Z$.
Therefore
\[
  D\cap Z\subseteq\bigcup_{n<N}(S_n\cap Z),
\]
a finite union of finite sets. Thus $|D\cap Z|<\omega$.

It follows that $D$ weakly separates $\mathcal Y$ from $\mathcal Z$, completing the proof.
\end{proof}


\begin{thm}\label{thm:at-equals-q2}
$\mathfrak{at}=\mathfrak q_{2}$.
\end{thm}

\begin{proof}
By Proposition~\ref{prop:q2-below-at}, it suffices to show that $\mathfrak{at}\leq \mathfrak q_{2}$.
Let $\kappa<\mathfrak{at}$.
We show that $\kappa<\mathfrak q_{2}$.
\end{proof}
